\chapter{一般相対性理論と多様体学習 --- 歪む「空間の形」が真理を教える}

\section*{本章の学習目標}
\begin{itemize}
\item ミンコフスキーがいかにしてユークリッド幾何学から「四次元時空」の数学を生み出したかを知る。
\item ミンコフスキー空間における「光円錐」と、宇宙の「因果律」の幾何学的な関係を理解する。
\item アインシュタインの恩師への反発と、その後の「幾何学への降伏」というドラマを学ぶ。
\item アインシュタインの「等価原理（重力と加速度は同じである）」という直感を理解する。
\item リーマン幾何学と測地線の概念から、「重力とは時空の歪みである」という真理を学ぶ。
\item 高次元データ空間に潜む「多様体仮説（Manifold Hypothesis）」の概念を理解する。
\item ディープラーニング（深層学習）の本質が、グニャグニャに曲がった空間の「アイロンがけ」であることを学ぶ。
\item オートエンコーダを例に、AIがいかにして曲がった空間の「本当の次元（座標）」を発見しているかを知る。
\end{itemize}

\section{光円錐と因果律：ミンコフスキーが描いた平らな時空}

前章で学んだ特殊相対性理論によって、時間と空間は独立した絶対的なものではなく、互いに混ざり合うものであることが明らかになりました。

アインシュタインが築いたこの新しい宇宙の舞台を、極めて美しく、かつ厳密な幾何学として描き直したのが、ドイツの天才数学者ヘルマン・ミンコフスキーです。

\subsection{ピタゴラスの定理から「四次元の幾何学」へ}

ミンコフスキーはもともと純粋数学（数論など）の世界的権威であり、計算式を美しい図形として捉える「数の幾何学」の創始者でした。彼は、かつての教え子であるアインシュタインが書いた特殊相対性理論の論文を読んだとき、その数式（ローレンツ変換）の背後に隠された「真の数学的構造」を瞬時に見抜きました。

私たちが学校で習うユークリッド幾何学では、空間の2点間の距離\(s\)は「ピタゴラスの定理（\(s^2=x^2+y^2+z^2\)）」で計算されます。

ミンコフスキーは、ここに「時間（\(t\)）」を4つ目の次元として付け加えようと考えました。しかし、時間を空間と同じようにそのまま足すわけにはいきません。彼は数学的な直感から、時間に光の速さ\(c\)と\textbf{「虚数\(i\)」}を掛け合わせた\(ict\)を、第4の座標として導入するという魔法を使いました。

虚数は2乗するとマイナス（\(-1\)）になるため、新しい4次元の距離は次のような数式になります。

\[
s^2=x^2+y^2+z^2-(ct)^2
\]

ピタゴラスの定理に、たった一つ「マイナス記号」がついた項が加わる。このわずかな違いが、通常のユークリッド幾何学とは全く異なる\textbf{「ミンコフスキー幾何学（擬ユークリッド幾何学）」}という新しい数学の分野を誕生させました。

このマイナス記号は、時間が単なる「4つ目の空間」ではないことを示しています。空間方向には自由に行ったり来たりできますが、時間方向には因果律という向きがあります。ミンコフスキーの式は、時間と空間を一つに統合しながらも、両者の性格の違いをきちんと残しているのです。

この四次元の距離（時空の間隔\(s\)）は、観測者がどれだけ猛スピードで動いて時間や空間が歪んで見えようとも、\textbf{誰から見ても絶対に変わらない「宇宙の不変量」}となることを彼は数学的に完全に証明したのです。

1908年、ケルンで開催された科学会議の伝説的な講演において、ミンコフスキーはこう高らかに宣言しました。

「今後、空間のみ、あるいは時間のみという概念は完全に影を潜め、両者が融合した『時空』のみが独立した実在となるだろう」

\subsection{「今、ここ」を中心とした光の限界}

私たちが「今、ここ（現在）」にいるとします。縦軸に「時間」、横軸に「空間（距離）」をとったグラフを想像してください。

光は宇宙で最も速く進む（秒速約30万km）ため、このミンコフスキーのグラフ上では一定の傾きを持った真っ直ぐな斜め線として描かれます。空間は実際には前後左右（3次元）に広がるため、この光の軌跡は「今、ここ」を頂点として、未来に向かってラッパ状に広がる\textbf{「円錐（コーン）」の形になります。同時に、過去から「今、ここ」に向かって収束してくる円錐も描けます。これが光円錐（Light Cone）}です。

\subsection{因果律の絶対的な境界線}

この光円錐は、単なる光の通り道ではなく、宇宙の絶対的なルールである\textbf{「因果律（原因と結果のつながり）」の境界線}を明確に引いています。

光円錐の内側（未来と過去）：光の速さ以下のスピードで移動したり、情報を伝えたりできる領域です。「今、ここ」であなたが投げたボールが届く範囲（確実な未来）や、「今、ここ」のあなたに影響を与えられた過去の出来事（確実な過去）は、すべてこの円錐の「内側」にすっぽりと収まります。

光円錐の外側（絶対的な他所：Elsewhere）：光のスピードを超えなければ到達できない領域です。相対性理論では光速を超えることはできないため、この円錐の外側にある出来事は、現在の私たちとは絶対に何の影響も及ぼし合うことができません。因果関係が完全に切り離された、物理的にアクセス不可能な「別の世界」です。

ミンコフスキー空間において、この光円錐の形（斜め線の角度）は、どんなに猛スピードで動いている観測者から見ても決して変わることはありません。宇宙の因果律は、平らで真っ直ぐな「光円錐」という幾何学的なバリアによって、美しく守られていたのです。

\section{アインシュタインの葛藤：恩師への反発と悲劇の別れ}

実は、この見事な「ミンコフスキー空間」が1908年に発表されたとき、アインシュタイン自身はこれを激しく嫌悪し、拒絶していました。

\subsection{「怠け者の犬」と「余計な虚飾」}

ミンコフスキーはチューリッヒ工科大学時代のアインシュタインの「数学の恩師」でした。しかし、当時のアインシュタインは物理学の直感ばかりを重んじ、数学の授業をサボってばかりいたため、ミンコフスキーからは\textbf{「あの怠け者の犬」と呆れられていたという因縁の師弟関係にありました。 アインシュタインは、自分の物理的な直感（特殊相対性理論）を、恩師が四次元の幾何学という極めて抽象的で難解な数学モデルに変換したのを見て、「あんなものは余計な数学的虚飾（衒学的な無用の長物）に過ぎない。数学者たちがいじくり回したせいで、私自身でさえ自分の理論が理解できなくなってしまった」}と毒づき、一切見向きもしなかったのです。

\subsection{ミンコフスキーの急死と残された呪縛}

しかし、運命は非情でした。あの伝説的な「時空」の講演からわずか数ヶ月後の1909年初頭、ミンコフスキーは盲腸炎（虫垂炎）をこじらせ、44歳という若さで突如この世を去ってしまいます。「彼がもう少し長く生きていれば、その後の物理学の歴史は大きく変わっていただろう」と惜しまれるほどの早すぎる死でした。

数年後、アインシュタインはこの自身の傲慢さを深く後悔することになります。

特殊相対性理論には「等速直線運動（スピードが一定で揺れない状態）」でしか使えないという致命的な弱点があり、彼はスピードが変化する「加速度」の世界と、ニュートンの最後の呪縛である\textbf{「重力（万有引力）」}の謎を解き明かそうと試みていました。

重力は、離れた星と星の間に瞬時に伝わる「遠隔作用」だと考えられていました。しかし、それでは前節で見た「光円錐のバリア（因果律）」を情報が一瞬で飛び越えてしまうことになり、相対性理論と真っ向から矛盾します。重力もまた、光速の制限（光円錐の中）に収めなければなりません。

アインシュタインは持ち前の「物理的直感」で重力に挑みましたが、どうやっても数式が辻褄の合わない矛盾だらけになってしまいます。

彼はついに気がつきました。「平らな時空」をいくらこねくり回しても重力は説明できない。重力を説明するには、\textbf{「時空そのものがグニャグニャに曲がっている」}という全く新しい数学的モデルが必要不可欠なのだ、と。

そして、その「曲がった時空」を記述するための絶対的な土台こそが、かつて彼自身が「無用の長物」と見下して切り捨てた、\textbf{亡き恩師が遺したミンコフスキー空間（四次元幾何学）}だったのです。

\subsection{「頼む、助けてくれ！」：数学へのSOS}

自分の数学力だけでは、平らなミンコフスキー幾何学を「曲がった空間を扱うリーマン幾何学（擬リーマン幾何学）」へと拡張・変形させることができないことに絶望したアインシュタインは、学生時代のノートを貸してくれていた親友であり、優秀な数学者になっていたマルセル・グロスマンに泣きつきました。

「グロスマン、頼む、私を助けてくれ！ さもないと私は気が狂ってしまう！」

かくしてアインシュタインは自らのプライドを捨て去り、物理の直感だけでなく「幾何学」という最強の数学的武器に完全に降伏し、それを受け入れました。後年、彼は「ミンコフスキーの幾何学的な枠組みがなければ、一般相対性理論は決して完成しなかっただろう」と深い敬意とともに述懐しています。

\section{人生最高のひらめき：「等価原理」と曲がる光}

グロスマンの助けを借りて高度な数学（テンソル解析やリーマン幾何学）を手に入れたアインシュタインは、1907年に得ていた「人生で最も素晴らしいひらめき」を、ついに厳密な理論へと昇華させていきます。

そのひらめきとは、\textbf{「屋根から落ちている最中の人には、自分の重さが感じられないだろう」}という事実でした。

これを宇宙空間の思考実験で考えてみましょう。

窓のない密室（エレベーター）の中にあなたが閉じ込められているとします。

地球上に置かれている場合：あなたは下に向かって「\(1G\)」の重力を感じ、床に足が押し付けられます。手からボールを離せば床に落ちます。

無重力の宇宙空間で、ロケットで上に向かって猛加速（\(1G\)の加速度）している場合：あなたは慣性によって床に押し付けられ、手からボールを離すと、床の方がボールに向かって加速してくるため、まるで「下に向かって重力があるように」感じます。

アインシュタインは、\textbf{「この2つの状況は、どんな物理実験を行っても絶対に区別できない。つまり、見かけの『加速度』と、星が引っ張る『重力』は、全く同じものの別の姿である」と見抜きました。これを「等価原理」}と呼びます。

この等価原理を認めると、驚くべき現象が予言されます。

無重力空間で猛加速しているエレベーターの横の小さな穴から、一筋の「光（レーザー）」が入ってきたとします。光がエレベーターを横切るほんのわずかな時間の間にも、エレベーター自体はどんどん上へと加速していくため、中にいる人から見ると、光の軌跡は真っ直ぐではなく「下に向かって放物線を描いて落ちた（曲がった）」ように見えます。

加速度と重力が同じであるならば、\textbf{「光は重力によって曲がって落ちる」}はずです。

しかし、ここで最大の矛盾が生じます。前章で見たように、\(E=mc^2\)の完全版方程式によれば「光には静止質量（重さ）がない」のです。ニュートンの万有引力（\(F=G\frac{mM}{r^2}\)）は、質量を持たないものには働きません。

質量ゼロの光が、なぜ重力に引っ張られて曲がるのでしょうか？

\section{一般相対性理論：重力とは「時空の幾何学」である}

アインシュタインが約10年の歳月とグロスマンから教わった幾何学を駆使してたどり着いた結論は、物理学の歴史を完全に書き換える壮大なものでした。重力は「力」ではなく、「曲がった空間の幾何学」だったのです。

\subsection{非ユークリッド幾何学の衝撃：平行線は交わる}

これを理解するために、私たちが学校で習う「ユークリッド幾何学（平らな紙の上の数学）」の常識を一度捨てなければなりません。

地球儀のような「曲がった表面」を想像してください。赤道上の離れた2点から、それぞれ真北に向かって「完全に平行な2本の直線」を引いて歩き出します。平らな紙の上なら平行線は永遠に交わりませんが、地球儀の上では、この2人は北極点で必ずドスンとぶつかって（交わって）しまいます。また、地球儀の上に三角形を描くと、内角の和は180度よりも大きくなります。

曲がった空間（非ユークリッド幾何学・リーマン幾何学の世界）においては、ピタゴラスの定理も、平行線のルールも成り立ちません。一見すると目に見えない引力で引き寄せられたように見えても、実は\textbf{「お互いが曲がった空間の表面に沿って真っ直ぐ進んだ結果、自然と距離が縮まってしまった」}だけなのです。

\subsection{トランポリンの宇宙：物質が空間を歪める}

アインシュタインは、この曲がった空間の数学を四次元の「時空」全体に適用しました。

「質量（エネルギー）を持った物質が存在すると、その周囲の『四次元時空』というミンコフスキーのキャンバスが、幾何学的に歪んでへこんでしまう。これが重力の正体である」

巨大なトランポリンの真ん中に、重いボウリングの球（太陽）を置いた状態を想像してください。球の重みで、シートはすり鉢状に深く沈み込み、歪みます。

その歪んだシートの端から、ビー玉（地球）を横に向かって転がしてみます。ビー玉自身は、曲がった空間の中で最も無駄のない真っ直ぐなルート（これを測地線：Geodesicと呼びます）を進もうとしているだけです。しかし、空間のシートそのものがすり鉢状に曲がっているため、外から見ると、結果としてボウリングの球の周りをクルクルと円を描いて回ることになります。

ニュートンには、このビー玉が「太陽からの目に見えない引力で引っ張られている」ように見えました。しかしアインシュタインは、\textbf{「太陽が空間を曲げ、地球はただその『曲がった空間の測地線』に沿って真っ直ぐ進んでいるだけだ」}と見破ったのです。

ただし、トランポリンの比喩には限界があります。トランポリンがへこむのは地球の重力が下向きに引っ張っているからですが、一般相対性理論で曲がるのは「下にある何か」ではなく、時空そのものです。また、本当に重要なのは空間だけでなく時間の曲がりです。惑星の軌道や光の曲がりは、四次元時空全体の幾何学として理解する必要があります。

\subsection{光の軌道と傾く光円錐：ブラックホールの誕生}

この幾何学的な視点に立てば、質量のない光が曲がる理由も完璧に説明できます。光は質量がないため引力には引かれませんが、空間そのものが曲がっている以上、光はその曲がった空間の形（測地線）に沿って進むしかないのです。だから光も曲がります。

これを7.1節で見た「光円錐」の視点で考えてみましょう。平らな宇宙では、光円錐は真っ直ぐ上（未来）に向かって立っていました。しかし、重力が強い星の近くでは、時空が歪むため、\textbf{「光円錐そのものが星の中心に向かってグニャリと傾いて（曲がって）しまう」}のです。

星が極限まで重くなり、ブラックホールのような凄まじい重力になると、この光円錐は完全に星の内側へと倒れ込んでしまいます。すると、光円錐の中にある「確実な未来」のすべてのルートが、星の中心（特異点）を向いてしまいます。光すらも因果律の外（外の世界）へ脱出することができなくなり、宇宙から切り離された絶対的な暗黒の穴となるのです。

\subsection{物質と時空の美しい対話}

物理学者ジョン・アーチボルド・ホイーラーは、この一般相対性理論の美しい関係性を次の一言で見事に表現しました。

「物質は、空間にどう曲がるかを教える。空間は、物質にどう動くかを教える」

宇宙とは、舞台（空間）と役者（星）が分かれた演劇ではありませんでした。役者の重みによって舞台そのものがグニャグニャと形を変え、その舞台の形によって役者の動きが決まり続ける。極めてダイナミックで相互作用的な「幾何学の織物」だったのです。

\section{AIが見つめる「曲がった宇宙」：多様体仮説}

「重力とは、引力ではなく空間の曲がり（幾何学）である」。

アインシュタインが「幾何学」への降伏を通じて到達したこの一般相対性理論の視点は、現代のAI（ディープラーニング）が複雑なデータを理解するための決定的な鍵を握っています。

ここでも、物理の時空とAIのデータ空間が同じものだという意味ではありません。共通しているのは、\textbf{「見えない力や意味を、空間の形として捉え直す」}という発想です。重力を力ではなく時空の曲がりとして見るように、AIは意味を単語や画像の孤立したラベルではなく、データ空間の曲がりや近さとして扱います。

前章で、AIが言葉を数百次元の「潜在空間（ベクトル）」に配置することを見ました。しかし、AIが扱う膨大なデータ空間は、ミンコフスキー空間のような真っ直ぐで平坦な空間（ユークリッド空間）ではありません。それは一般相対論の重力場のように、極めて複雑に歪み、グニャグニャに曲がりくねっています。

\subsection{1万次元の砂嵐と「薄いシート」}

例えば、100×100ピクセルの画像データは、1万個の数字の集まりなので数学的には「1万次元空間」の1点になります。

しかし、この1万次元空間にデタラメに点を打っても、ただのテレビの砂嵐のようなノイズ画像にしかなりません。現実の「犬の顔」や「手書きの数字」といった「意味のあるデータ」は、この広大な1万次元空間のあらゆる場所に均等に散らばっているわけではないのです。

なぜなら、現実の画像には\textbf{「強い物理的な制約」}があるからです。「隣り合うピクセルは似た色になりやすい」「光は上から当たる」「犬には目が2つある」といった物理世界のルールや幾何学的な構造が、見えない重力のように働いて、取り得るデータのパターンを極限まで絞り込んでいます。

この少数の「真のルール（顔の向き、光の強さ、犬種のパラメータなど）」こそが、1万次元の虚無の空間の中に浮かび上がる、\textbf{極めて薄く、かつ複雑に折りたたまれた「低次元のシート」の正体です。 この曲がったシートのことを数学用語で「多様体（Manifold）」と呼び、現実の複雑なデータは低次元の多様体上にのみ存在するという考え方を「多様体仮説（Manifold Hypothesis）」}と呼びます。

分かりやすく言えば、私たちが住む地球の表面（2次元の多様体）が、宇宙空間（3次元）の中で丸く曲がって存在しているのと同じことです。

\subsection{直線距離が通用しない世界：AIにとっての「測地線」}

従来のAI（古い機械学習）が画像認識などを苦手としていた最大の理由は、この「空間の曲がり」を理解できず、直線定規（ユークリッド距離）で無理やり距離を測ろうとしたからです。

例えば、「少し右を向いた犬（画像A）」から「少し左を向いた犬（画像B）」へと画像を滑らかに変化（モーフィング）させたいとします。

もし1万次元の平らな空間上で、画像Aから画像Bへと「真っ直ぐな直線」を引いてその中間のピクセルを計算（平均化）するとどうなるでしょうか？ 結果は、右を向いた犬と左を向いた犬が半透明に重なり合った、不気味な「顔が2つあるお化けのような犬」のノイズ画像になってしまいます。多様体（意味のあるデータのシート）から完全に外れて、砂嵐が広がる虚無の空間に飛び出してしまうからです。

本当に正しい「中間の画像（正面を向いた犬）」を見つけるためには、空間に平らな直線を引くのではなく、グニャグニャに曲がった多様体のシートの表面にピタリと張り付いたまま、AからBへの最短ルートを辿らなければなりません。

地球上の東京からニューヨークへの本当の距離が、地球の内部を突き抜ける直線ではなく、曲がった地表に沿ったルートで測らなければならないのと同じです。

お気づきでしょうか。この「曲がった空間の表面に沿った最短ルート」こそ、まさにアインシュタインが重力の中で地球や光が進む軌道として見破った\textbf{「測地線（Geodesic）」}に他なりません。

一般相対性理論において、物質が「曲がった時空のシート上の測地線」に沿って進むように、現代の生成AIもまた、意味を破綻させることなくデータとデータの間を繋ぐために「曲がったデータ空間上の測地線」を計算して進まなければならないのです。

\section{ディープラーニングの本質：空間の「アイロンがけ」}

では、最新のAI（ディープラーニング）は、このグニャグニャに曲がった多様体の形をどのようにして見破り、「本当の距離（意味）」を測れるようになったのでしょうか。

ディープラーニング（多層ニューラルネットワーク）がやっている究極の仕事の一つは、\textbf{「高次元空間の中で複雑に絡み合った多様体（データのシート）を、層を経るごとに少しずつ引っ張り、綺麗に平らに引き伸ばすこと（空間のアイロンがけ）」}として理解できます。

もちろん、すべてのニューラルネットワークが文字通り多様体を完全に平坦化しているわけではありません。しかし、複雑に絡んだデータを、分類や生成がしやすい表現へと少しずつ変換していくという意味では、この「アイロンがけ」の比喩はディープラーニングの直感をよく捉えています。

\subsection{折り紙とハサミ：行列演算と活性化関数}

ニューラルネットワークの1つの層では、主に2つの数学的な操作が交互に行われています。

空間を伸ばす・回す（アフィン変換）：データに「重み（行列）」を掛け算することで、空間全体を斜めに引き伸ばしたり、回転させたりします。

空間を折る・切る（活性化関数）：引き伸ばしただけでは、空間の「曲がり具合（非線形性）」は変えられません。そこで「ReLU」と呼ばれるような活性化関数を通します。これは、「マイナスの数値になったら強制的にゼロにする（切り捨てる）」という乱暴な操作です。この操作によって、空間のシートに「折り目」をつけたり、不要な部分を「ハサミで切り取ったり」することができるようになります。

AIのネットワークは、この\textbf{「行列で引っ張り、活性化関数で折り目をつけたり切り取ったりする」という折り紙のような操作}をひたすら繰り返すことで、複雑に曲がった多様体を強引に平らな形へと変形させていくのです。

\subsection{スイスロールの解きほぐしと「深さ（Deep）」の理由}

この「アイロンがけ」のプロセスを想像するために、機械学習の世界でよく使われる\textbf{「スイスロール」}というデータ構造を思い浮かべてください。これは、2次元の平らなシートが、3次元空間の中でロールケーキのようにぐるぐると渦を巻いて巻き取られている状態です。

もしこれを「たった1回の操作（1つの層）」で無理やり真っ直ぐに引き伸ばそうとすると、シートがビリビリに破れてしまい、データの意味（繋がり）が完全に壊れてしまいます。

渦を巻いたシートを破らずに綺麗に平らにするためには、「少し引き出しては伸ばし、また少し引き出しては伸ばし\ldots{}\ldots{}」というように、小さな変形を何段階にも分けて慎重に行う必要があります。

これこそが、ディープラーニングが\textbf{「深く（Deepに）」ならなければならない物理的・幾何学的な理由}です。

第1層がデータの表面的な浅いシワを伸ばし、第2層がより深い折り目を伸ばし、第10層、第100層と進むにつれて、スイスロールのように複雑に丸まったデータ空間の歪みを少しずつ解きほぐし、最終的に「一枚の平らなシート（人間が理解できる特徴量の空間）」へと完璧にアイロンがけしていくのです。

\subsection{オートエンコーダ：意味の重力場を発見する}

この多様体のアイロンがけのプロセスを最も純粋な形で見せてくれるのが、\textbf{「オートエンコーダ（自己符号化器）」}と呼ばれるAI技術です。

オートエンコーダは、1万次元の入力データを、ネットワークの途中でわざと「100次元」のような極端に狭いボトルネック（潜在空間）に押し込み、そこから再び元の1万次元のデータを復元しようとするAIです。

これはAIにとって極めて過酷な試練です。情報を極限まで圧縮されるため、AIは「画像にとって本当に重要な本質」だけを選び出さなければなりません。この学習を通じて、AIは1万次元の空間の大部分を占める「意味のない虚無の空間（ノイズ）」を切り捨て、データが実際に張り付いている「多様体そのものの形」だけを正確に学習します。

それはまるで、一般相対性理論において、グニャグニャに曲がった時空の表面にピッタリと沿うような「新しい目盛り（計量テンソル）」を見つけ出す作業そのものです。

学習を終えたオートエンコーダのボトルネック（潜在空間）の中には、複雑な多様体が綺麗にアイロンがけされ、真っ直ぐな定規が使えるようになった\textbf{「平らで美しい意味の宇宙」}が広がっています。

前章で見た「King - Man + Woman = Queen」のようなベクトルの直線計算ができるようになったのは、AIがディープラーニングという強力なアイロンを使って、重力場のように歪んでいたデータ空間のシートを見事に「平坦化」したからなのです。

\section{まとめ：形こそが真理である}

アインシュタインは自分の直感だけでは限界があることを悟り、かつて否定した恩師の「幾何学」へと降伏することで、重力という見えない力の正体を「空間の曲がり」へと変換し、宇宙の真の姿を解き明かしました。

現代のAIもまた、「犬らしさや猫らしさという見えない意味」の正体を、「高次元空間における多様体の曲がりという幾何学」へと変換することで、データに潜む真理を解き明かしています。

人類が外に広がるマクロな宇宙の形を捉えたその全く同じ数学と視点（幾何学）を用いて、AIは今、情報というミクロな宇宙の形を捉えようとしているのです。

\section*{【第7章 章末ディスカッション課題】}

光円錐の外側は「絶対に情報が伝わらない（因果律の外）」と学びました。もし仮に、光の速さを超えて情報を送れるタイムマシン（タキオン通信）が発明されたとしたら、原因と結果の順番（因果律）はどう崩壊してしまうでしょうか？

もし私たちが「1万次元空間」の曲がり方を直感的に想像できる脳を持っていたとしたら、ディープラーニングのようなAIの助けを借りなくても、複雑な画像や言語の意味を一瞬で理解できるでしょうか？ 人間の脳の幾何学的な限界とAIの役割について考えてみましょう。


\section*{参考文献（学術誌）}
\begin{enumerate}[leftmargin=2.4em]
\item Einstein, A. ``Die Grundlage der allgemeinen Relativitatstheorie.'' \textit{Annalen der Physik}, 49, 769--822, 1916. DOI: \url{https://doi.org/10.1002/andp.19163540702}.
\item Dyson, F. W., Eddington, A. S., and Davidson, C. ``A determination of the deflection of light by the Sun's gravitational field, from observations made at the total eclipse of May 29, 1919.'' \textit{Philosophical Transactions of the Royal Society A}, 220, 291--333, 1920. DOI: \url{https://doi.org/10.1098/rsta.1920.0009}.
\item Tenenbaum, J. B., de Silva, V., and Langford, J. C. ``A global geometric framework for nonlinear dimensionality reduction.'' \textit{Science}, 290(5500), 2319--2323, 2000. DOI: \url{https://doi.org/10.1126/science.290.5500.2319}.
\item Belkin, M. and Niyogi, P. ``Laplacian eigenmaps for dimensionality reduction and data representation.'' \textit{Neural Computation}, 15(6), 1373--1396, 2003. DOI: \url{https://doi.org/10.1162/089976603321780317}.
\end{enumerate}


